\section{支集\index{zhiji@支集|hyperpage}与 $n$-同调}
\label{sec:03.3.9}

本节要证明若干关于 $n$-同调的结果，它们来自对交织函子作用方式的分析。

先作一些几何准备。

设 $S$ 是旗簇 $X$ 的子集。对 $w\in W$，置
$$E_w(S)=\{x\in X\mid \bb_x\ \text{相对于}\ \bb_y\ \text{处于相对位置}\ v,\ \text{其中}\ v\le w,\ y\in S\}.$$

\begin{lemma}\label{lem:03.9.1}
\begin{enumerate}[label=(\roman*),leftmargin=2.6em,itemsep=0.15em]
\item 设 $S\subset X$，$w\in W$。则
$$\dim S\le\dim E_w(S)\le\dim S+\ell(w).$$
\item 设 $S\subset X$，$w\in W$。则
$$E_w(\overline S)=E_w(S).$$
\item 设 $S$ 是 $X$ 的闭子集，$w\in W$。则 $E_w(S)$ 是集合
$$\{x\in X\mid \bb_x\ \text{相对于某个}\ \bb_y\ (y\in S)\ \text{处于相对位置}\ w\}$$
的闭包。
\item 若 $S$ 不可约，则 $E_w(S)$ 也不可约。
\item 设 $w,v\in W$ 满足 $\ell(wv)=\ell(w)+\ell(v)$。则
$$E_{wv}(S)=E_w(E_v(S)).$$
\end{enumerate}
\end{lemma}

\begin{proof}
设 $\alpha\in\Pi$。记 $X_\alpha$ 为全体 $\alpha$ 型抛物子代数构成的簇，$p_\alpha:X\to X_\alpha$ 为自然投影。则
\begin{align*}
E_{s_\alpha}(S)&=\{x\in X\mid \bb_x\ \text{相对于}\ \bb_y\ \text{处于相对位置}\ v,\ v\le s_\alpha,\ y\in S\}\\
&=S\cup\{x\in X\mid \bb_x\ \text{相对于某个}\ \bb_y\ (y\in S)\ \text{处于相对位置}\ s_\alpha\}=p_\alpha^{-1}(p_\alpha(S)).
\end{align*}
显然，因 $p_\alpha:X\to X_\alpha$ 是纤维同构于 $\mathbb{P}^1$ 的局部平凡纤维化，故当 $S$ 闭（相应地，不可约）时 $E_{s_\alpha}(S)$ 也闭（相应地，不可约）。此外还有
$$\dim S\le\dim E_{s_\alpha}(S)\le\dim S+1.$$
因此 $E_{s_\alpha}(\overline S)$ 闭，故 $E_{s_\alpha}(S)\subset E_{s_\alpha}(\overline S)$。另一方面，由 $S\subset E_{s_\alpha}(S)$ 可得 $\overline S\subset E_{s_\alpha}(S)$。若 $x\in E_{s_\alpha}(S)$，则整条纤维 $p_\alpha^{-1}(p_\alpha(x))$ 都含于 $E_{s_\alpha}(S)$。这说明 $E_{s_\alpha}(\overline S)\subset E_{s_\alpha}(S)$。这就对单反射证明了 (ii)。

现在对 $w\in W$ 的长度作归纳来证 (v)。先断言 $w=s_\alpha$（$\alpha\in\Pi$）时公式成立。此时 $E_{s_\alpha}(E_v(S))$ 由这样的 $x\in X$ 组成：或者 $x\in E_v(S)$，或者存在 $y\in E_v(S)$ 使 $\bb_x$ 相对于 $\bb_y$ 处于相对位置 $s_\alpha$。因此它由这样的 $x\in X$ 组成：存在 $y\in S$，使 $\bb_x$ 相对于 $\bb_y$ 处于相对位置 $u$，其中或者 $u\le v$，或者 $u=s_\alpha u'$ 且 $u'\le v$。在第二种情形，或者 $\ell(u)=\ell(u')+1$ 且 $u\le s_\alpha v$，或者 $\ell(u)=\ell(u')-1$ 且 $u\le u'\le v$。故 $E_{s_\alpha}(E_v(S))\subset E_{s_\alpha v}(S)$。反之，若 $u\le s_\alpha v$，则或者 $u\le v$，或者 $s_\alpha u\le v$，故 $E_{s_\alpha}(E_v(S))=E_{s_\alpha v}(S)$。

现设 $w$ 任意。可取 $\alpha\in\Pi$ 与 $w'\in W$，使 $w=s_\alpha w'$ 且 $\ell(w)=\ell(w')+1$。于是由归纳假设，
$$E_w(E_v(S))=E_{s_\alpha w'}(E_v(S))=E_{s_\alpha}(E_{w'}(E_v(S)))=E_{s_\alpha}(E_{w'v}(S)),$$
这就完成了 (v) 的证明。

现在设 $w\in W$、$\alpha\in\Pi$、$w'\in W$ 任意，且 $w=s_\alpha w'$、$\ell(w)=\ell(w')+1$。此时 $E_w(S)=E_{s_\alpha}(E_{w'}(S))$。用证明的第一部分并对 $\ell(w)$ 作归纳，(i)、(ii)、(iv) 立即得出。此外还看出：若 $S$ 闭，则 $E_w(S)$ 闭。

(iii) 置
$$V=\{x\in X\mid \bb_x\ \text{相对于某个}\ \bb_y\ (y\in S)\ \text{处于相对位置}\ w\}.$$
则 $V\subset E_w(S)$。因 $E_w(S)$ 闭，故 $\overline V\subset E_w(S)$。取 $y\in S$。则全体满足「$\bb_x$ 相对于 $\bb_y$ 处于相对位置 $w$」的 $x\in X$ 所成集合的闭包等于 $E_w(\{y\})$。这说明
$$\overline V\supset\bigcup_{y\in S}E_w(\{y\})=E_w(S).$$
\end{proof}

称 $w\in W$ 与 $S\subset X$ \emph{横截}，若
$$\dim E_w(S)=\dim S+\ell(w).$$
若 $w$ 与 $S$ 横截，则 $\ell(w)\le\operatorname{codim} S$。

\begin{lemma}\label{lem:03.9.2}
\begin{enumerate}[label=(\roman*),leftmargin=2.6em,itemsep=0.1em]
\item $w\in W$ 与 $S$ 横截，当且仅当它与 $\overline S$ 横截。
\item 设 $S\subset X$，$w,v\in W$ 满足 $\ell(wv)=\ell(w)+\ell(v)$。则下列陈述等价：
\begin{enumerate}[label=(\alph*)]
\item $wv$ 与 $S$ 横截；
\item $v$ 与 $S$ 横截，且 $w$ 与 $E_v(S)$ 横截。
\end{enumerate}
\end{enumerate}
\end{lemma}

\begin{proof}
(i) 由 9.1(ii) 有 $\dim E_w(S)=\dim E_w(\overline S)$，断言由横截性的定义得出。

(ii) 由 9.1(i) 有
$$\dim E_{wv}(S)\le\dim S+\ell(wv)=\dim S+\ell(w)+\ell(v),$$
且等号成立当且仅当 $wv$ 与 $S$ 横截。另一方面，由 9.1(v) 有
$$\dim E_{wv}(S)=\dim E_w(E_v(S))\le\dim E_v(S)+\ell(w)\le\dim S+\ell(v)+\ell(w).$$
于是若 (a) 成立，上式中的不等号全取等号，即
$$\dim E_w(E_v(S))=\dim E_v(S)+\ell(w),\qquad \dim E_v(S)=\dim S+\ell(v).$$
故 (b) 成立。反之，若 (b) 成立，则立即看出 $wv$ 与 $S$ 横截。
\end{proof}

\begin{lemma}\label{lem:03.9.3}
设 $S$ 是 $X$ 的不可约闭子簇，$w\in W$。则存在 $v\le w$，使 $v$ 与 $S$ 横截且 $E_v(S)=E_w(S)$。
\end{lemma}

\begin{proof}
先考虑 $w=s_\alpha$（$\alpha\in\Pi$）的情形。此时 $E_{s_\alpha}(S)=p_\alpha^{-1}(p_\alpha(S))$ 不可约且闭，且只有两种可能：
(a) $s_\alpha$ 与 $S$ 横截，且 $\dim E_{s_\alpha}(S)=\dim S+1$；
(b) $s_\alpha$ 与 $S$ 不横截，$\dim E_{s_\alpha}(S)=\dim S$，且由 $S\subset E_{s_\alpha}(S)$ 得 $E_{s_\alpha}(S)=S$。

现对 $\ell(w)$ 作归纳来证一般情形。若 $\ell(w)=0$，则 $w=1$ 而 $E_1(S)=S$，断言显然。设 $\ell(w)=k$。则存在 $w'\in W$ 与 $\alpha\in\Pi$，使 $w=s_\alpha w'$ 且 $\ell(w)=\ell(w')+1$。此时由 9.1(v) 有 $E_w(S)=E_{s_\alpha}(E_{w'}(S))$。由归纳假设，存在 $v'\le w'$，它与 $S$ 横截，且 $E_{v'}(S)=E_{w'}(S)$。

现在，由证明的第一部分，若 $s_\alpha$ 与 $E_{w'}(S)$ 不横截，则
$$E_w(S)=E_{s_\alpha}(E_{w'}(S))=E_{w'}(S)=E_{v'}(S).$$
因 $v'\le w'\le w$，断言成立。若 $s_\alpha$ 与 $E_{w'}(S)$ 横截，则
$$\dim E_w(S)=\dim E_{s_\alpha}(E_{w'}(S))=\dim E_{w'}(S)+1=\dim S+\ell(v')+1.$$
置 $v=s_\alpha v'$。若 $\ell(v)=\ell(v')-1$，则由 9.1(v) 有
$$E_{v'}(S)=E_{s_\alpha}(E_v(S))=p_\alpha^{-1}\bigl(p_\alpha(E_v(S))\bigr),$$
从而
$$E_{s_\alpha}(E_{v'}(S))=p_\alpha^{-1}\bigl(p_\alpha(p_\alpha^{-1}(p_\alpha(E_v(S))))\bigr)=E_{v'}(S),$$
这与 $s_\alpha$ 与 $E_{w'}(S)$ 的横截性矛盾。故 $\ell(v)=\ell(v')+1$，$v\le w$，且 $E_v(S)=E_{s_\alpha}(E_{v'}(S))$。于是 $E_w(S)=E_v(S)$，且
$$\dim E_v(S)=\dim E_w(S)=\dim S+\ell(v')+1=\dim S+\ell(v),$$
即 $v$ 与 $S$ 横截。
\end{proof}

固定 $\lambda\in\h^*$。记 $\mathcal{M}_{coh}(\D_\lambda)$ 为凝聚 $\D_\lambda$-模构成的范畴。凝聚 $\D_\lambda$-模 $\mathcal{V}$ 的支集 $\operatorname{supp}\mathcal{V}$ 是 $X$ 的闭子簇~[9]。我们要分析交织函子的作用如何改变凝聚 $\D$-模的支集。

先指出一个简单事实，它是交织函子定义与 9.1(iii) 的直接推论。设 $V^\bullet$ 是 $\D^b(\D_\lambda)$ 中的复形，定义其支集为
$$\operatorname{supp}V^\bullet=\bigcup_{p\in\Z}\operatorname{supp}H^p(V^\bullet).$$
显然，由上述注记，$V^\bullet\in\D^b_{coh}(\D_\lambda)$ 的支集是 $X$ 的闭子簇。

\begin{lemma}\label{lem:03.9.4}
对任意 $V^\bullet\in\D^b_{coh}(\D_\lambda)$ 与 $w\in W$ 有
$$\operatorname{supp}LI_w(V^\bullet)\subset E_w(\operatorname{supp}V^\bullet).$$
\end{lemma}

\begin{proof}
先对单反射证明这一结果。设 $\alpha\in\Pi$。则
$$LI_{s_\alpha}(V^\bullet)=Rp_{1+}\bigl(T_{s_\alpha}\otimes_{\calO_{Z_{s_\alpha}}}p_2^+(V^\bullet)\bigr),$$
且有谱序列
$$R^sp_{1+}\bigl(T_{s_\alpha}\otimes_{\calO_{Z_{s_\alpha}}}H^t(p_2^+(V^\bullet))\bigr)\Longrightarrow H^{s+t}(LI_{s_\alpha}(V^\bullet)).$$
故 $LI_{s_\alpha}(V^\bullet)$ 的支集含于 $p_2^+(V^\bullet)$ 的支集之像的闭包中。而 $p_2^+(V^\bullet)$ 的支集含于 $Z_{s_\alpha}$ 的闭子集
{\small
\begin{align*}
&\{(x,x')\in Z_{s_\alpha}\mid x'\in\operatorname{supp}V^\bullet\}\\
&\qquad=\{(x,x')\in X\times X\mid \bb_x\ \text{相对于}\ \bb_{x'}\ \text{处于相对位置}\ s_\alpha,\ x'\in\operatorname{supp}V^\bullet\}
\end{align*}}
中。此集合在 $p_1$ 下的投影等于
$$\{x\in X\mid \bb_x\ \text{相对于}\ \bb_{x'}\ \text{处于相对位置}\ s_\alpha,\ x'\in\operatorname{supp}V^\bullet\}\subset E_{s_\alpha}(\operatorname{supp}V^\bullet).$$
故 $\operatorname{supp}LI_{s_\alpha}(V^\bullet)\subset E_{s_\alpha}(\operatorname{supp}V^\bullet)$。

现对 $\ell(w)$ 作归纳来证一般情形。设 $w=w's_\alpha$，其中 $w'\in W$，$\alpha\in\Pi$，且 $\ell(w)=\ell(w')+1$。则由 3.18 有 $LI_w=LI_{w'}\circ LI_{s_\alpha}$。于是由归纳假设与 9.1(v)，
\begin{align*}
\operatorname{supp}LI_w(V^\bullet)&=\operatorname{supp}LI_{w'}(LI_{s_\alpha}(V^\bullet))\subset E_{w'}(\operatorname{supp}LI_{s_\alpha}(V^\bullet))\\
&\subset E_{w'}(E_{s_\alpha}(\operatorname{supp}V^\bullet))=E_w(\operatorname{supp}V^\bullet).
\end{align*}
\end{proof}

下一个结果更精细。

\begin{lemma}\label{lem:03.9.5}
设 $\mathcal{V}\in\mathcal{M}_{coh}(\D_\lambda)$，$\alpha\in\Pi$。则
\begin{enumerate}[label=(\roman*)]
\item $\dim\operatorname{supp}L^{-1}I_{s_\alpha}(\mathcal{V})\le\dim\operatorname{supp}\mathcal{V}$；
\item 若 $s_\alpha$ 与 $\operatorname{supp}\mathcal{V}$ 横截，则 $\dim\operatorname{supp}I_{s_\alpha}(\mathcal{V})=\dim\operatorname{supp}\mathcal{V}+1$。
\end{enumerate}
\end{lemma}

\begin{proof}
置 $S=\operatorname{supp}\mathcal{V}$。则只有两种可能：
(a) $s_\alpha$ 与 $S$ 横截；
(b) $s_\alpha$ 与 $S$ 不横截。

先看情形 (b)。此时 $\dim S=\dim E_{s_\alpha}(S)$。于是由 9.4，对 $p\in\Z$ 有
$$\dim\operatorname{supp}L^pI_{s_\alpha}(\mathcal{V})\le\dim S.$$
特别地，此时 (i) 成立。

剩下要研究情形 (a)。先作一些几何准备。设 $S_i$（$1\le i\le n$）是 $S$ 的不可约分支。则
$$E_{s_\alpha}(S)=E_{s_\alpha}\Bigl(\bigcup_{i=1}^n S_i\Bigr)=\bigcup_{i=1}^n E_{s_\alpha}(S_i).$$
由 9.1(iv)，$E_{s_\alpha}(S_i)$ 都闭且不可约，故族 $(E_{s_\alpha}(S_i);1\le i\le n)$ 的极大元就是 $E_{s_\alpha}(S)$ 的不可约分支。于是
$$\dim E_{s_\alpha}(S)=\max_{1\le i\le n}\dim E_{s_\alpha}(S_i),$$
且存在 $S$ 的不可约分支 $S_i$ 使 $\dim E_{s_\alpha}(S_i)=\dim E_{s_\alpha}(S)$。重排下标后，可设这对 $1\le i\le m$ 成立。因 $s_\alpha$ 与 $S$ 横截，有 $\dim E_{s_\alpha}(S)=\dim S+1$。故当 $1\le i\le m$ 时 $\dim E_{s_\alpha}(S_i)=\dim S+1$。另一方面，$\dim E_{s_\alpha}(S_i)\le\dim S_i+1$ 说明 $1\le i\le m$ 时 $\dim S_i=\dim S$。于是 $S_i$（$1\le i\le m$）都是 $S$ 的、维数为 $\dim S$ 的不可约分支。因 $E_{s_\alpha}(S_i)=p_\alpha^{-1}(p_\alpha(S_i))$，而 $p_\alpha:X\to X_\alpha$ 是纤维同构于 $\mathbb{P}^1$ 的局部平凡纤维化，故当 $1\le i\le m$ 时
$$\dim p_\alpha(S_i)=\dim p_\alpha^{-1}(p_\alpha(S_i))-1=\dim S_i=\dim S.$$
另一方面，若 $m<i\le n$，则 $\dim E_{s_\alpha}(S_i)\le\dim S$，故或者 $\dim S_i<\dim S$，或者 $s_\alpha$ 与 $S_i$ 不横截。两种情形下都有 $\dim p_\alpha(S_i)<\dim S$。

于是，若记 $S_0$ 为 $S$ 的奇点集与 $\bigcup_{m<i\le n}S_i$ 的并，则 $\dim p_\alpha(S_0)<\dim S$，且 $\dim E_{s_\alpha}(S_0)\le\dim S$。置 $X'=X-S_0$。则 $S\cap X'$ 是 $X'$ 的光滑闭子簇，其不可约分支为 $S'_i=S_i\cap X'$（$1\le i\le m$）。于是 $S'_i$（$1\le i\le m$）是维数为 $\dim S$ 的、互不相交的光滑子簇，且 $s_\alpha$ 与它们都横截。

若考虑限制 $p_\alpha|_{S'_i}:S'_i\to p_\alpha(S'_i)$（$1\le i\le m$），则存在 $p_\alpha(S_i)$ 中的开稠密集 $U_i$，使 $u\in U=\bigcup_{i=1}^m U_i$ 时纤维 $p_\alpha^{-1}(u)\cap S$ 有限~[11, Ch.~I, \S8, Theorem~3]。集合 $p_\alpha^{-1}(U)$ 在 $E_{s_\alpha}(S)$ 中开，维数为 $\dim S+1$。因 $S$ 与 $E_{s_\alpha}(S_0)$ 都是 $E_{s_\alpha}(S)$ 的余维 $1$ 的闭子空间，故
$$V=p_\alpha^{-1}(U)-(S\cup E_{s_\alpha}(S_0))$$
是 $E_{s_\alpha}(S)$ 中维数为 $\dim S+1$ 的开集，其补集是维数为 $\dim S$ 的子簇。

取 $x\in V$。考虑投影 $p_i:Z_{s_\alpha}\to X$（$i=1,2$），它们由 $X\times X$ 到第一、第二因子的投影诱导。则
\begin{align*}
p_1^{-1}(x)\cap p_2^{-1}(S)&=\{(x,x')\in X\times X\mid p_\alpha(x)=p_\alpha(x'),\ x'\in S\}\\
&=\{(x,x')\mid x'\in p_\alpha^{-1}(p_\alpha(x))\cap S\}
\end{align*}
是非空有限集，且 $p_2$ 给出此集合到 $p_\alpha^{-1}(p_\alpha(x))\cap S$ 的双射。

现在来分析 $LI_{s_\alpha}$ 的几何纤维。对任意 $x\in X$，记 $i_x$ 为自然单射 $\{x\}\to X$。则对任意 $U^\bullet\in\D^b(\D_\lambda)$ 有
$$LT_x(U^\bullet)=Li_x^+(U^\bullet)=Ri_x^!(U^\bullet)[\dim X].$$
因此对任意 $x\in X$ 有
\begin{align*}
LT_x(LI_{s_\alpha}(D(\mathcal{V})))&=Ri_x^!\bigl(LI_{s_\alpha}(D(\mathcal{V}))\bigr)[\dim X]\\
&=Ri_x^!\Bigl(Rp_{1+}\bigl(T_{s_\alpha}\otimes_{\calO_{Z_{s_\alpha}}}p_2^+(D(\mathcal{V}))\bigr)\Bigr)[\dim X].
\end{align*}
设 $Z_x=p_1^{-1}(x)\subset Z_{s_\alpha}$ 是 $p_1$ 在 $x$ 上的纤维，$j_x$ 是 $Z_x$ 到 $Z_{s_\alpha}$ 的浸入。则由基变换~[1, VI.8.4] 有
{\small
\begin{align*}
LT_x(LI_{s_\alpha}(D(\mathcal{V})))&=Ri_x^!\Bigl(Rp_{1+}\bigl(T_{s_\alpha}\otimes_{\calO_{Z_{s_\alpha}}}p_2^+(D(\mathcal{V}))\bigr)\Bigr)[\dim X]\\
&=R(p_1\circ j_x)_+\Bigl(Rj_x^!\bigl(T_{s_\alpha}\otimes_{\calO_{Z_{s_\alpha}}}p_2^+(D(\mathcal{V}))\bigr)\Bigr)[\dim X]\\
&=R(p_1\circ j_x)_+\bigl(j_x^*(T_{s_\alpha})\otimes_{\calO_{Z_x}}Rj_x^!(p_2^+(D(\mathcal{V}))[-1])\bigr)[\dim X]\\
&=R(p_1\circ j_x)_+\bigl(j_x^*(T_{s_\alpha})\otimes_{\calO_{Z_x}}R(p_2\circ j_x)^!(D(\mathcal{V}))\bigr)[\dim X-1].
\end{align*}}
投影 $p_2\circ j_x$ 是 $Z_x$ 到其像 $F=p_\alpha^{-1}(p_\alpha(x))-\{x\}\subset X$ 的同构。记 $q_1:F\to X$ 与 $q_2:F\to X$ 分别为 $p_1$ 与 $p_2$ 与此同构之逆的复合。则 $q_2$ 是 $F$ 到 $X$ 的自然包含。因由 3.3(i) 有 $T_{s_\alpha}=p_1^*(\calO(\alpha))=p_2^*(\calO(-\alpha))$，最终得到
$$LT_x(LI_{s_\alpha}(D(\mathcal{V})))=Rq_{1+}\bigl(q_2^*(\calO(-\alpha))\otimes_{\calO_F}Rq_2^!(D(\mathcal{V}))\bigr)[\dim X-1].$$

现设 $x\in V$。则 $F\subset X'$。另一方面，$S'=S\cap X'$ 是 $X'$ 的光滑闭子簇。设 $j$ 是 $S'$ 到 $X$ 的自然浸入。则由 Kashiwara 的范畴等价~[1, VI.7.11] 有
$$\mathcal{V}|_{X'}=j_+\bigl(R^0j^!(\mathcal{V})\bigr)|_{X'}.$$
故
$$Rq_2^!(D(\mathcal{V}))=Rq_2^!\bigl(D(j_+(R^0j^!(\mathcal{V})))\bigr)=Rq_2^!\bigl(j_+(Rj^!(D(\mathcal{V})))\bigr).$$
设 $k$ 与 $h$ 分别是 $F\cap S$ 到 $F$ 与 $S'$ 的自然浸入。再用一次基变换，得到
\begin{align*}
Rq_2^!(D(\mathcal{V}))&=Rq_2^!\bigl(j_+(Rj^!(D(\mathcal{V})))\bigr)=Rk_+\bigl(Rh^!(Rj^!(D(\mathcal{V})))\bigr)\\
&=Rk_+\bigl(R(j\circ h)^!(D(\mathcal{V}))\bigr)=Rk_+\bigl(R(q_2\circ k)^!(D(\mathcal{V}))\bigr).
\end{align*}
因 $F\cap S$ 有限，若记 $k_y$ 为 $\{y\}$ 到 $F$ 的自然浸入，则
\begin{align*}
Rq_2^!(D(\mathcal{V}))&=Rk_+\bigl(R(q_2\circ k)^!(D(\mathcal{V}))\bigr)=\bigoplus_{y\in F\cap S}k_{y+}\bigl(Ri_y^!(D(\mathcal{V}))\bigr)\\
&=\bigoplus_{y\in F\cap S}k_{y+}\bigl(LT_y(D(\mathcal{V}))\bigr)[-\dim X].
\end{align*}
这说明
$$LT_x(LI_{s_\alpha}(D(\mathcal{V})))=\bigoplus_{y\in(p_\alpha^{-1}(p_\alpha(x))-\{x\})\cap S}LT_y(D(\mathcal{V}))[-1]$$
作为向量空间复形成立。

由~[1, V.9.3]，再把 $U$ 缩小一些，可设 $j^!(\mathcal{V})$ 是 $S\cap p_\alpha^{-1}(U)$ 上的局部自由 $\calO$-模。这说明当 $p\ne 0$、$y\in S\cap p_\alpha^{-1}(U)$ 时 $L^pT_y(j^!(\mathcal{V}))=0$；且当 $y\in S\cap p_\alpha^{-1}(U)$ 时
$$E_y=T_y\bigl(j^!(\mathcal{V})\bigr)\ne 0.$$
因此，若记 $l_y$ 为 $\{y\}$ 到 $S'$ 的包含，则对所有 $y\in S\cap p_\alpha^{-1}(U)$ 有
{\small
\begin{align*}
LT_y(D(\mathcal{V}))&=Ri_y^!(D(\mathcal{V}))[\dim X]=Rl_y^!\bigl(Rj^!(D(\mathcal{V}))\bigr)[\dim X]\\
&=LT_y\bigl(D(j^!(\mathcal{V}))\bigr)[\dim X-\dim S]=D(E_y)[\dim X-\dim S].
\end{align*}}
于是对 $x\in V$ 有
{\small
$$LT_x(LI_{s_\alpha}(D(\mathcal{V})))=\bigoplus_{y\in(p_\alpha^{-1}(p_\alpha(x))-\{x\})\cap S}D(E_y)[\dim X-\dim S-1],$$}
且 $x\in\operatorname{supp}LI_{s_\alpha}(D(\mathcal{V}))$。因此 $\operatorname{supp}LI_{s_\alpha}(D(\mathcal{V}))$ 含 $V$ 的闭包 $\overline V$，故
$$\dim\operatorname{supp}LI_{s_\alpha}(D(\mathcal{V}))\ge\dim S+1.$$
而由 9.4 有 $\dim\operatorname{supp}LI_{s_\alpha}(D(\mathcal{V}))=\dim S+1$。
因
\begin{align*}
\dim\operatorname{supp}LI_{s_\alpha}(D(\mathcal{V}))&=\max_{p\in\Z}\dim\operatorname{supp}L^pI_{s_\alpha}(\mathcal{V})\\
&=\max\bigl\{\dim\operatorname{supp}I_{s_\alpha}(\mathcal{V}),\ \dim\operatorname{supp}L^{-1}I_{s_\alpha}(\mathcal{V})\bigr\},
\end{align*}
为完成证明，还需说明 $\dim\operatorname{supp}L^{-1}I_{s_\alpha}(\mathcal{V})\le\dim S$。设 $X''$ 是 $E_{s_\alpha}(S)$ 的奇点集与其维数不超过 $\dim S$ 的不可约分支之并的补集。则 $T=E_{s_\alpha}(S)\cap X''$ 是 $X''$ 的闭光滑子簇，其所有不可约分支的维数都是 $\dim S+1$。设 $v:T\to X$ 是自然包含。由 Kashiwara 定理，$L^pI_{s_\alpha}(\mathcal{V})|_{X''}=v_+(C_p)|_{X''}$，其中 $C_p=v^!(L^pI_{s_\alpha}(\mathcal{V}))$ 是 $T$ 上的凝聚 $\D$-模。必要时缩小 $X''$，可设 $C_p$ 都是局部自由 $\calO$-模~[1, VII.9.3]。若记 $h_x$ 为 $\{x\}$ 到 $T$ 的包含，再用一次基变换，可见对 $x\in T$ 与 $p\in\Z$ 有
{\small
\begin{align*}
LT_x\bigl(D(L^pI_{s_\alpha}(\mathcal{V}))\bigr)&=LT_x\bigl(D(v_+(C_p))\bigr)=Ri_x^!\bigl(v_+(D(C_p))\bigr)[\dim X]\\
&=Rh_x^!\bigl(D(C_p)\bigr)[\dim X]=LT_x\bigl(D(C_p)\bigr)[\dim X-\dim S-1]\\
&=D\bigl(T_x(C_p)\bigr)[\dim X-\dim S-1].
\end{align*}}
因此当 $q\ne\dim S-\dim X+1$ 时
$$L^qT_x\bigl(L^pI_{s_\alpha}(\mathcal{V})\bigr)=0,$$
且
$$L^{\dim S-\dim X+1}T_x\bigl(L^pI_{s_\alpha}(\mathcal{V})\bigr)=T_x(C_p).$$
由这一计算，谱序列
$$L^qT_x\bigl(L^pI_{s_\alpha}(\mathcal{V})\bigr)\Longrightarrow H^{p+q}\bigl(LT_x(LI_{s_\alpha}(D(\mathcal{V})))\bigr)$$
在 $E_2$ 项退化，且对 $x\in T$ 有
$$H^{\dim S-\dim X+p+1}\bigl(LT_x(LI_{s_\alpha}(D(\mathcal{V}))))=T_x(C_p).$$
于是对 $x\in T$ 有
{\small
$$LT_x(LI_{s_\alpha}(D(\mathcal{V})))=D\bigl(T_x(C_{-1})\bigr)[\dim X-\dim S]\oplus D\bigl(T_x(C_0)\bigr)[\dim X-\dim S-1].$$}
因 $E_{s_\alpha}(S)$ 的不可约分支的集合等于 $(E_{s_\alpha}(S_i);1\le i\le m)$，可见 $T\cap V$ 在 $T$ 中稠密。比较前面的计算，可见当 $x\in T\cap V$ 时 $T_x(C_{-1})=0$。因 $C_p$ 局部自由，故 $C_{-1}=0$。于是 $\operatorname{supp}L^{-1}I_{s_\alpha}(\mathcal{V})\subset X-X''$，且 $\dim\operatorname{supp}L^{-1}I_{s_\alpha}(\mathcal{V})\le\dim S$。这就完成了情形 (a) 中 (i) 的证明，并说明 (ii) 必成立。
\end{proof}

\begin{proposition}\label{prop:03.9.6}
设 $\mathcal{V}\in\mathcal{M}_{coh}(\D_\lambda)$，$w\in W$。则对 $p\in\Z$ 有
$$\dim\operatorname{supp}L^pI_w(\mathcal{V})\le\dim\operatorname{supp}\mathcal{V}+\ell(w)+p.$$
\end{proposition}

\begin{proof}
对 $\ell(w)$ 作归纳。若 $w$ 是单反射，这由 9.1(i)、9.4 与 9.5(i) 得出。设 $w=s_\alpha w'$，其中 $\alpha\in\Pi$，$w'\in W$，且 $\ell(w)=\ell(w')+1$。则由归纳假设，对 $p,q\in\Z$ 有
{\small
\begin{align*}
\dim\operatorname{supp}L^pI_{w'}\bigl(L^qI_{s_\alpha}(\mathcal{V})\bigr)&\le\dim\operatorname{supp}L^qI_{s_\alpha}(\mathcal{V})+\ell(w')+p\\
&\le\dim\operatorname{supp}\mathcal{V}+\ell(w')+1+p+q=\dim\operatorname{supp}\mathcal{V}+\ell(w)+p+q.
\end{align*}}
由附着于 3.18 的谱序列可见，对任意 $s\in\Z$ 有
$$\dim\operatorname{supp}L^sI_w(\mathcal{V})\le\dim\operatorname{supp}\mathcal{V}+\ell(w)+s.$$
\end{proof}

\begin{lemma}\label{lem:03.9.7}
设 $\mathcal{V}\in\mathcal{M}_{coh}(\D_\lambda)$，$w\in W$ 与 $\operatorname{supp}\mathcal{V}$ 横截，且 $\mathcal{V}$ 的支集 $S$ 不可约。则
$$\operatorname{supp}I_w(\mathcal{V})=E_w(S).$$
\end{lemma}

\begin{proof}
对 $\ell(w)$ 作归纳。若 $\ell(w)=1$，则 $w=s_\alpha$（$\alpha\in\Pi$）。由 9.4，$\operatorname{supp}I_{s_\alpha}(\mathcal{V})\subset E_{s_\alpha}(S)$。又由 9.1，两个集合都闭，且 $E_{s_\alpha}(S)$ 不可约。由横截性与 9.5(ii) 有 $\dim\operatorname{supp}I_{s_\alpha}(\mathcal{V})=\dim S+1=\dim E_{s_\alpha}(S)$，断言随之成立。

设 $w\in W$，$\ell(w)=k>1$。则 $w=s_\alpha w'$，其中 $\alpha\in\Pi$，$\ell(w')=k-1$。因 $w$ 与 $S$ 横截，由 9.2 知 $w'$ 与 $S$ 横截，且 $s_\alpha$ 与 $E_{w'}(S)$ 横截。由归纳假设，$\operatorname{supp}I_{w'}(\mathcal{V})=E_{w'}(S)$。于是由 3.8 与 9.1(iv) 有
$$\operatorname{supp}I_w(\mathcal{V})=\operatorname{supp}I_{s_\alpha}\bigl(I_{w'}(\mathcal{V})\bigr)=E_{s_\alpha}(E_{w'}(S))=E_w(S).$$
\end{proof}

对每个凝聚 $\D_\lambda$-模，联系 $W$ 的两个子集：
$$S(\mathcal{V})=\{w\in W\mid\operatorname{supp}I_w(\mathcal{V})=X\},$$
以及
$$E(\mathcal{V})=S(\mathcal{V})\ \text{中极小元所成的集合}.$$

\begin{proposition}\label{prop:03.9.8}
设 $\mathcal{V}\in\mathcal{M}_{coh}(\D_\lambda)$，且其支集不可约。则
\begin{enumerate}[label=(\roman*)]
\item $S(\mathcal{V})$ 非空；
\item $E(\mathcal{V})=\{w\in W\mid w\ \text{与}\ \operatorname{supp}\mathcal{V}\ \text{横截，且}\ \ell(w)=\operatorname{codim}\operatorname{supp}\mathcal{V}\}$，即 $E(\mathcal{V})$ 由所有与 $\operatorname{supp}\mathcal{V}$ 横截、且长度达到最大可能的 $w\in W$ 组成。
\end{enumerate}
\end{proposition}

\begin{proof}
设 $w\in W$ 与 $\operatorname{supp}\mathcal{V}$ 横截，且 $\ell(w)=\operatorname{codim}\operatorname{supp}\mathcal{V}$。则由 9.7 得 $w\in S(\mathcal{V})$。若 $v<w$，则 $\ell(v)<\operatorname{codim}\operatorname{supp}\mathcal{V}$，且由 9.4 有 $\dim\operatorname{supp}I_v(\mathcal{V})<\dim X$。故 $v\notin S(\mathcal{V})$，即 $w\in E(\mathcal{V})$。

反之，设 $w\in E(\mathcal{V})$。则由 9.4 有 $E_w(\operatorname{supp}\mathcal{V})=X$。因 $\mathcal{V}$ 的支集不可约，由 9.3 可取 $v\le w$，使 $v$ 与 $\operatorname{supp}\mathcal{V}$ 横截，且 $E_v(\operatorname{supp}\mathcal{V})=X$。由 9.7，这说明 $v\in S(\mathcal{V})$。因 $w$ 是 $S(\mathcal{V})$ 中的极小元，必须 $w=v$，故 $w$ 与 $\operatorname{supp}\mathcal{V}$ 横截。这就证明了 (ii)。

为证 (i)，只需说明 $E(\mathcal{V})$ 非空。显然，若 $w_0$ 是 $W$ 中的最长元，则 $E_{w_0}(S)=X$。由 9.3，存在与 $S$ 横截的 $w$ 使 $E_w(S)=X$，故断言由 (ii) 得出。
\end{proof}

为陈述本节的主要结果，还需要一个概念。设 $V$ 是有限生成 $U_\theta$-模。称 $\lambda\in\theta$ 为 $V$ 的\emph{指数}，若集合
$$\{x\in X\mid H_0(\n_x,V)_{(\lambda+\rho)}\ne 0\}$$
含 $X$ 的一个开稠密子集。

称 $\lambda\in\h^*$ \emph{强反支配}，若对任意 $\alpha\in\Sigma^+$ 有 $\operatorname{Re}\check{\alpha}(\lambda)\le 0$。显然，强反支配的 $\lambda$ 反支配。

再在 $\h^*$ 上定义偏序：$\lambda\preccurlyeq\mu$，若 $\mu-\lambda$ 是 $\Pi$ 中单根的线性组合，且系数实部非负。这个序关系与 $W$ 上的序有如下联系。

\begin{lemma}\label{lem:03.9.9}
设 $\lambda\in\h^*$ 强反支配。则对任意 $v,w\in W$，$v\le w$ 推出 $v\lambda\preccurlyeq w\lambda$。
\end{lemma}

\begin{proof}
显然，只需对任意 $w\in W$ 与 $\alpha\in\Pi$ 证明：若 $\ell(s_\alpha w)=\ell(w)+1$，则 $w\lambda\preccurlyeq s_\alpha w\lambda$。但 $s_\alpha w\lambda=w\lambda-\check{\alpha}(w\lambda)\alpha$，故
$$s_\alpha w\lambda-w\lambda=\bigl(w^{-1}\alpha\bigr)^{\vee}(\lambda)\alpha,$$
于是只需证明 $\operatorname{Re}\bigl(w^{-1}\alpha\bigr)^{\vee}(\lambda)\ge 0$。因 $w^{-1}\alpha\in\Sigma^+$~[5, Ch.~VI, \S1, no.~6, Cor.~2 of Prop.~17]，这立即由 $\lambda$ 的强反支配性得出。
\end{proof}

\begin{theorem}\label{thm:03.9.10}
设 $\lambda\in\h^*$ 强反支配。设 $\mathcal{V}\in\mathcal{M}_{coh}(\D_\lambda)$，且 $S=\operatorname{supp}\mathcal{V}$ 不可约。置 $V=\Gamma(X,\mathcal{V})$。
\begin{enumerate}[label=(\roman*)]
\item 若 $\omega$ 是 $V$ 的指数，则存在与 $S$ 横截、且 $\ell(w)=\operatorname{codim}S$ 的 $w\in W$，使 $w\lambda\preccurlyeq\omega$。
\item 设 $V$ 不可约且 $V\ne 0$。若 $w\in W$ 与 $S$ 横截，且 $\ell(w)=\operatorname{codim}S$，则 $w\lambda$ 是 $V$ 的指数。
\end{enumerate}
\end{theorem}

\begin{proof}
(i) 设 $\mu$ 是正则支配权，$F$ 是以 $\mu$ 为最高权的不可约有限维 $\g$-模。置 $\mathcal{F}=\calO_X\otimes_\C F$。则 $\lambda-\mu$ 正则且强反支配。置 $U=\Gamma(X,\mathcal{V}(-\mu))$。则由 C.2.1 有
$$\mathcal{V}=\bigl(\mathcal{V}(-\mu)\otimes_{\calO_X}\mathcal{F}\bigr)_{[\lambda]}.$$
这说明
{\small
\begin{align*}
V&=\Gamma(X,\mathcal{V})=\Gamma\bigl(X,(\mathcal{V}(-\mu)\otimes_{\calO_X}\mathcal{F})_{[\lambda]}\bigr)\\
&=\bigl(\Gamma(X,\mathcal{V}(-\mu)\otimes_{\calO_X}\mathcal{F})\bigr)_{[\lambda]}
=\bigl(\Gamma(X,\mathcal{V}(-\mu))\otimes_\C F\bigr)_{[\lambda]}=(U\otimes_\C F)_{[\lambda]}.
\end{align*}}
设 $\omega$ 是 $V$ 的指数，即对 $X$ 的某个开稠密子集中的所有 $x$ 有 $H_0(\n_x,V)_{(\omega+\rho)}\ne 0$。则
$$H_0(\n_x,V)=H_0\bigl(\n_x,(U\otimes_\C F)_{[\lambda]}\bigr)$$
是 $H_0(\n_x,U\otimes_\C F)$ 中相应于权 $v\lambda+\rho$（$v\in W$）的广义 $U(\h)$-特征空间之直和。故
$$H_0(\n_x,V)_{(\omega+\rho)}=H_0(\n_x,U\otimes_\C F)_{(\omega+\rho)}.$$
设 $(F_p;1\le p\le n)$ 是 $F$ 中递增的 $\bb_x$-不变极大旗。它诱导 $\bb_x$-模 $U\otimes_\C F$ 的滤过 $(U\otimes_\C F_p;1\le p\le n)$。相应的分次模是形如 $U\otimes_\C\C_\nu$ 的模之和，其中 $\nu$ 取遍 $F$ 的所有权。显然，$H_0(\n_x,U\otimes_\C F)$ 的半单化是模 $H_0(\n_x,U)\otimes_\C\C_\nu$ 的直和的子模。因 $U$ 的无穷小特征标正则，由 L.2.4，$H_0(\n_x,U)$ 是半单 $\h$-模。这说明 $H_0(\n_x,V)_{(\omega+\rho)}$ 是模 $H_0(\n_x,U)_{(\omega-\nu+\rho)}\otimes_\C\C_\nu$ 的直和的子模。特别地，若 $H_0(\n_x,V)_{(\omega+\rho)}\ne 0$，则对 $F$ 的某个权 $\nu$ 有 $H_0(\n_x,U)_{(\omega-\nu+\rho)}\ne 0$。因权集有限，可设对 $X$ 的某个开稠密子集中的所有 $x$ 有 $H_0(\n_x,U)_{(\omega-\nu+\rho)}\ne 0$。另一方面，$\omega-\nu=v(\lambda-\mu)$，其中 $v\in W$ 由 $\omega-\nu$ 唯一确定。这说明 $v^{-1}(\omega-\nu)=\lambda-\mu$。因 $\omega=u\lambda$（$u\in W$），可见
$$v^{-1}u\lambda-\lambda=-(\mu-v^{-1}\nu).$$
因 $\mu$ 是 $F$ 的最高权，右端是正根之和的相反数。又因 $\lambda$ 反支配，左端是正根之和。于是两端都必须为零：$v^{-1}u$ 属于 $\lambda$ 的稳定化子，且 $\omega=u\lambda=v\lambda$。

因 $\lambda-\mu$ 正则，有 $\mathcal{V}(-\mu)=\Delta_{\lambda-\mu}(U)$。此外，由 2.6 可得 $\operatorname{supp}\Delta_{v(\lambda-\mu)}(U)=X$。因由 3.16 有
$$I_v(\mathcal{V}(-\mu))=I_v(\Delta_{\lambda-\mu}(U))=\Delta_{v(\lambda-\mu)}(U),$$
故 $v\in S(\mathcal{V}(-\mu))=S(\mathcal{V})$。于是由 9.3，存在 $w\le v$，它与 $S$ 横截，且 $\ell(w)=\operatorname{codim}S$。而由 9.9，这推出 $w\lambda\preccurlyeq v\lambda=\omega$。

(ii) 若 $V$ 不可约，则 $\mathcal{V}(-\mu)$ 也不可约，且二者的支集 $S$ 不可约。故由范畴等价，$U$ 不可约。因 $w$ 与 $S$ 横截且 $\ell(w)=\operatorname{codim}S$，由 9.7 可见 $\operatorname{supp}\Delta_{w(\lambda-\mu)}(U)=X$。置 $\mathcal{U}=\Delta_{w(\lambda-\mu)}(U)$。因 $U$ 不可约，对 $U$ 用 1.16（取 $p=0$）即得 $U\subset\Gamma(X,\mathcal{U})$。

设 $s\in U$ 是 $\mathcal{U}$ 的整体截面，且它在 $X$ 的某个开稠密子集上为零。于是它生成一个支集含于该开集补集中的整体截面子模。这个子模或者等于 $U$，或者等于零。第一种可能说明局部化 $\Delta_{w(\lambda-\mu)}(U)$ 的支集也含于该开集的补集中，与假设矛盾。故这个子模等于零，即 $s=0$。这说明 $\mathcal{U}$ 中任何非零整体截面的支集都等于 $X$。

设 $F$ 是以 $\mu$ 为最高权的 $\g$ 的不可约有限维表示。则如前，由 C.2.1 有
$$\mathcal{U}(w\mu)=\bigl(\mathcal{U}\otimes_{\calO_X}\mathcal{F}\bigr)_{[\lambda]}.$$
于是
{\small
\begin{align*}
\Gamma(X,\mathcal{U}(w\mu))&=\Gamma\bigl(X,(\mathcal{U}\otimes_{\calO_X}\mathcal{F})_{[\lambda]}\bigr)=\bigl(\Gamma(X,\mathcal{U}\otimes_{\calO_X}\mathcal{F})\bigr)_{[\lambda]}\\
&=\bigl(\Gamma(X,\mathcal{U})\otimes_\C F\bigr)_{[\lambda]}\supset(U\otimes_\C F)_{[\lambda]}=V.
\end{align*}}
此外，$\mathcal{U}\otimes_{\calO_X}\mathcal{F}=\mathcal{U}\otimes_\C F$ 的任何非零整体截面，只要它来自 $U\otimes_\C F$，其支集都等于 $X$；其子层 $\mathcal{U}(w\mu)$ 中属于 $(U\otimes_\C F)_{[\lambda]}=V$ 的任何非零整体截面，支集也等于 $X$。因 $\mathcal{U}(w\mu)$ 凝聚，存在 $X$ 中的开稠密集 $O$，使 $\mathcal{U}(w\mu)|_O$ 是局部自由 $\calO_O$-模~[1, VII.9.3]。于是在这个集合上，截面为零当且仅当其值（即其在几何纤维中的像）处处为零。故存在 $O$ 的开稠密子集 $O'$，使对 $x\in O'$，来自 $V$ 的某些截面在 $x$ 处不为零。另一方面，对任意 $x\in O'$，$\n_xV$ 中的整体截面都在该点为零。因此，对 $x\in O'$，几何纤维映射 $\mathcal{U}(w\mu)\mapsto T_x(\mathcal{U}(w\mu))$ 诱导出 $V$ 到 $T_x(\mathcal{U}(w\mu))$ 的非零映射，它经 $H_0(\n_x,V)$ 分解，且这个分解映射是 $\bb_x$-模态射。由此可见，对 $x\in O'$ 有 $H_0(\n_x,V)_{(w\lambda+\rho)}\ne 0$，即 $w\lambda$ 是 $V$ 的指数。
\end{proof}

\begin{annote}
底本此处有两处编号与上下文不合。其一，「于是由 6.」处要用的存在性（存在 $w\le v$，与 $S$ 横截且 $\ell(w)=\operatorname{codim}S$）出自 9.3；9.6 只是 $\dim\operatorname{supp}L^pI_w$ 的估计，给不出这样的 $w$。其二，「而由 7.」处，由 $w\le v$ 推出 $w\lambda\preccurlyeq v\lambda$ 是 9.9 的内容；9.7 只给出支集，不含 $\lambda$ 的序关系。已按上下文改为 9.3 与 9.9。[来源：底本 p.83，9.10(i) 的证明；对照 9.3、9.6、9.7、9.9 的陈述]
\end{annote}

下一个结果是 9.10 的直接推论。

\begin{theorem}\label{thm:03.9.11}
设 $V\ne 0$ 是有限生成 $U_\theta$-模。则 $V$ 的指数集非空。特别地，存在 $X$ 的开稠密子集 $U$，使对 $x\in U$ 有 $H_0(\n_x,V)\ne 0$。
\end{theorem}

\begin{proof}
因 $V$ 非零，它有不可约商 $U\ne 0$。取强反支配的 $\lambda\in\theta$。则由 4.2，存在不可约 $\D_\lambda$-模 $\mathcal{U}$ 使 $U=\Gamma(X,\mathcal{U})$。由 9.8(ii)，存在 $\omega\in\theta$，它是 $U$ 的指数。因 $H_0(\n_x,U)$ 是 $H_0(\n_x,V)$ 的商，故 $\omega$ 也是 $V$ 的指数。
\end{proof}

\begin{annote}
底本末句作「故 $\omega$ 是 $U$ 的指数」。$U$ 是 $V$ 的不可约商，$\omega$ 先由 9.10(ii) 对 $U$ 得到；定理要证的是 $V$ 的指数集非空，而正因为 $H_0(\n_x,U)$ 是 $H_0(\n_x,V)$ 的商，$\omega$ 才成为 $V$ 的指数。故此处应为 $V$，已改。[来源：底本 p.84，9.11 的陈述与证明]
\end{annote}
